\documentclass[11pt, fleqn]{amsart}
\usepackage{amsmath, amssymb, amsthm}
\usepackage{hyperref}
\usepackage{mathtools}
\usepackage{listings}
\usepackage{xcolor}
\lstset{
  basicstyle=\ttfamily\small,
  breaklines=true,
  frame=single,
  language=Python,
  keywordstyle=\color{blue},
  commentstyle=\color{gray},
  stringstyle=\color{red}
}
\newtheorem{theorem}{Theorem}[section]
\newtheorem{lemma}[theorem]{Lemma}
\newtheorem{corollary}[theorem]{Corollary}
\newtheorem{proposition}[theorem]{Proposition}
\theoremstyle{remark}
\newtheorem{remark}{Remark}

\title{Spectral graph theory and the matrices $J_n(\overline{j})$, $H_n(\overline{h})$:
remarks on the digraph viewpoint}
\date{12 June 2026}

\begin{document}
\maketitle

\noindent
These remarks concern the two matrices defined after Lemma 2.1 of
\emph{Oscillating spectra associated to certain sequences}, regarded as
weighted adjacency matrices of directed graphs. The short answer is yes:
spectral graph theory says quite a bit, and the directed-graph viewpoint
fits these matrices unusually well. The remarks are organized from
``structural'' to ``potentially useful for the $\mu_n$ project.''
References to proofs of every result invoked are given inline, and
Section~\ref{sec:guide} is an annotated guide to the bibliography
intended for a reader with little background in spectral graph theory.

\section{The digraph of $J_n$ has a very rigid cycle structure}

Read $J_n(\overline{j})$ as the weighted adjacency matrix of a digraph on
vertices $1,\dots,n$, with the convention that the entry $A_{uv}$ is the
weight of the edge $u\to v$. (For the general dictionary between square
matrices and weighted digraphs, see Brualdi and Cvetkovi\'c
\cite{BC09}, the gentlest introduction we know, or Chapter~6 of Horn
and Johnson \cite{HJ13}.) The superdiagonal gives a single
``ascending'' path $i \to i+1$ with weight $-i$, and the lower-triangular
Toeplitz part gives ``descending'' edges $i \to k$ ($k \le i$) with weight
$j_{i-k+1}$, including loops of weight $j_1$. Because the \emph{only} way
to ascend is one step at a time, every directed cycle is supported on an
interval $[k,m]$: it climbs $k \to k+1 \to \cdots \to m$ and returns via
the single edge $m \to k$. Its weight is
\begin{equation}\label{eq:cycleweight}
w([k,m]) \;=\; j_{m-k+1}\,(-1)^{m-k}\,\frac{(m-1)!}{(k-1)!}.
\end{equation}

\section{The coefficient theorem for digraphs reproves Theorem 2.3
combinatorially}

The \emph{coefficient theorem} for weighted digraphs --- in the
graph-theory literature it is usually attributed to Sachs and appears,
stated and proved for digraphs, in Chapter~1 of Cvetkovi\'c, Doob, and
Sachs \cite{CDS95} --- says that the coefficient of $x^{n-r}$ in
$\chi(J_n)$ is
\[
\sum_{L} (-1)^{c(L)}\, w(L),
\]
the sum running over collections $L$ of vertex-disjoint cycles covering
$r$ vertices, where $c(L)$ is the number of cycles in $L$ and $w(L)$ is
the product of the edge weights occurring in $L$. (The theorem is an
unpacking of the Leibniz expansion of $\det(xI-A)$ via the cycle
decomposition of permutations. Self-contained proofs accessible without
graph-theoretic background may be found in Brualdi--Cvetkovi\'c
\cite{BC09}, where determinants and characteristic polynomials are
developed from the start in the language of weighted digraphs, and in
Stanley \cite[\S 4.7]{Stanley12}, as the first step of the
transfer-matrix method.)

By the previous section, the relevant collections $L$ are exactly sets
of pairwise disjoint intervals in $\{1,\dots,n\}$ --- that is, the
coefficients are signed sums over partial composition-like structures,
which is morally why partition and composition identities (MacDonald,
Vein--Dale) appear at all in this setting.

This was verified numerically (see the script in
Section~\ref{sec:script}): for a random integer sequence $\overline{j}$
with $j_1=1$ and $n=7$, the cycle-cover sums match both \texttt{numpy}'s
characteristic polynomial and the formula of Theorem 2.3,
\[
\binom{n}{r}\, r!\, h_r\, (-1)^r,
\]
exactly.

A corollary worth noting: for digraphs the spectrum depends \emph{only}
on the multiset of cycle products. One direction is immediate from the
coefficient theorem, since the coefficients of $\chi$ are assembled
from cycle products alone. The deeper converse-type statement --- two
matrices with the same digraph are diagonally similar if and only if
all corresponding cycle products agree --- is a theorem of Saunders and
Schneider \cite{SS78}; diagonal similarity is the matrix form of a
``gauge transformation'' that reweights edges while preserving cycle
products. Hence the interval weights \eqref{eq:cycleweight} are the
complete intrinsic data behind the spectra under study. Any answer to
the question ``what is it about $\chi(J_n)$ that induces the observed
oscillations?'' must be expressible in terms of them.

\section{Strong connectivity and its standard dividends}

The digraph is strongly connected (one can ascend stepwise and descend
in one hop), so $J_n$ is irreducible: the equivalence
``irreducible $\Leftrightarrow$ strongly connected digraph'' is proved
in \cite[Ch.~6]{HJ13} and in \cite[Ch.~2]{BP94}. Moreover, since the
superdiagonal is nowhere zero, $J_n$ is an \emph{unreduced} Hessenberg
matrix, hence nonderogatory --- every eigenvalue has geometric
multiplicity $1$, and $\chi(J_n)$ coincides with the minimal
polynomial. (Proof in one line: deleting the first column and last row
of $J_n - \lambda I$ leaves an upper-triangular matrix with diagonal
$-1, -2, \dots, -(n-1)$, which is nonsingular, so
$\operatorname{rank}(J_n - \lambda I) \ge n-1$ for every $\lambda$. The
standard treatment of unreduced Hessenberg matrices is Golub and Van
Loan \cite[\S 7.4]{GVL13}.)

Irreducibility upgrades the Ger\v{s}gorin disc theorem
\cite{Ger31} via Taussky's theorem: an eigenvalue lying on the boundary
of the union of the discs must lie on the boundary of \emph{every}
disc. Taussky's original article is \cite{Taussky49}; modern proofs of
both the disc theorem and Taussky's refinement appear in
\cite[Ch.~6]{HJ13} and, with much more leisurely exposition, in
Chapter~1 of Varga's book \cite{Varga04}, which is devoted entirely to
this circle of ideas and is written to be readable without prior
background. More interestingly, Brualdi's cycle-based refinement of
Ger\v{s}gorin is spectral graph theory tailor-made for the present
situation: it localizes eigenvalues using products of deleted row sums
taken \emph{around directed cycles} of the digraph, and here the cycles
are exactly the intervals described above. Brualdi's theorem, with
proof, is in his original paper \cite{Brualdi82}; a textbook proof,
together with comparisons against the classical discs and Brauer's
ovals of Cassini, is in Chapter~2 of \cite{Varga04}.

\section{The Ger\v{s}gorin picture explains, geometrically, what
deformation does}

Every diagonal entry of $J_n$ equals $j_1$, so all Ger\v{s}gorin discs
share one center. Deformation by $c$ replaces every loop weight by $c$
(and shifts the descending-edge weights), so it translates the common
disc center to $c$ while reorganizing the radii. In graph language,
deformation is a uniform reweighting of the loops --- the spectrum is
pinned to circulate around $c$. This gives at least a heuristic for why
the deformed spectra look more regular: the eigenvalues live in an
annulus-like region about a fixed center, and $\mu_n$ measures the
distance from $0$ to a cloud orbiting $c$.

\section{A concrete lever for bounding $\mu_n$ away from zero}

Since the determinant is the product of the eigenvalues (with
multiplicity),
\[
\prod_i |\lambda_i| \;=\; |\det J_n| \;=\; n!\,|h_n|,
\]
any upper bound $\rho$ on the spectral radius coming from the graph
(Ger\v{s}gorin row sums in terms of the $j$'s --- see
\cite[Ch.~6]{HJ13} or \cite[Ch.~1]{Varga04} --- or Brualdi's cycle
bounds \cite{Brualdi82}) yields
\[
\mu_n \;\ge\; \frac{n!\,|h_n|}{\rho^{\,n-1}}.
\]
This is exactly the shape of statement the project wants --- it reduces
``the spectrum is zero-free'' to controlling the top of the spectrum
graph-theoretically, plus knowing $h_n \neq 0$. For the Lehmer
application the logic unfortunately runs the wrong way ($h_n \neq 0$ is
the desired conclusion, not a hypothesis). Still, for the deformed
matrices, where only zero-freeness is needed, it is potentially usable.

\section{Caveats, and the matrix $H_n$}

One important caveat: the classical heavy machinery of spectral graph
theory assumes nonnegative or symmetric matrices. The
Perron--Frobenius theorem (proved in \cite[Ch.~8]{HJ13}, in
\cite[Ch.~2]{BP94}, and --- in what is perhaps the most readable
account --- in Chapter~8 of Meyer \cite{Meyer00}) requires nonnegative
entries; eigenvalue interlacing and its many graph-theoretic
consequences (see \cite[\S 4.3]{HJ13} for the matrix statement, or
Brouwer and Haemers \cite{BH12} for the spectral-graph-theory
applications) require symmetry. The weights here are signed and the
matrices are far from normal, so the spectra are genuinely complex
point sets and those tools do not apply. What survives is exactly the
digraph toolkit above: the coefficient theorem, irreducibility, and
cycle-based inclusion regions.

Everything said for $J_n$ transfers to $H_n(\overline{h})$ with the
obvious changes: the superdiagonal weights are all $+1$, and the
first-column edges $i \to 1$ carry weight $i\,h_i$. Consequently the
cycles of the $H_n$ digraph are the loops, the intervals $[k,m]$ with
$k \ge 2$, and the intervals $[1,m]$, the last carrying the extra
factor $m$.

\section{Guide to the references}\label{sec:guide}

For a reader new to spectral graph theory we suggest the following
order. References below are given at the chapter or section level
rather than by theorem number, since numbering differs between
editions.

\emph{Start here.} Brualdi--Cvetkovi\'c \cite{BC09} develops elementary
matrix theory (determinants, characteristic polynomials,
Cayley--Hamilton) entirely through weighted digraphs and cycle covers,
assuming only linear algebra; it contains everything needed for
Sections 1--2 of these remarks, with proofs. Stanley
\cite[\S 4.7]{Stanley12} gives a compact alternative proof of the
cycle-cover expansion as part of the transfer-matrix method.

\emph{The eigenvalue-inclusion world.} Varga \cite{Varga04} is a short,
self-contained book devoted to Ger\v{s}gorin-type theorems: Chapter~1
proves the disc theorem and Taussky's boundary refinement, and
Chapter~2 proves Brualdi's cycle-based inclusion sets used above. The
original sources are Ger\v{s}gorin \cite{Ger31}, Taussky
\cite{Taussky49}, and Brualdi \cite{Brualdi82}; Brualdi's paper is also
the best place to see how the digraph cycle structure enters the
proofs.

\emph{Standard references.} Horn--Johnson \cite{HJ13} (already cited in
the main paper for Proposition 2.4 there) contains proofs of every
matrix-theoretic fact we used: Ger\v{s}gorin, Taussky, and
irreducibility in Chapter~6, interlacing in \S 4.3, Perron--Frobenius
in Chapter~8. Golub--Van Loan \cite[\S 7.4]{GVL13} covers unreduced
Hessenberg matrices. Berman--Plemmons \cite{BP94} and Meyer
\cite{Meyer00} are alternatives for Perron--Frobenius and
irreducibility, Meyer being the gentlest.

\emph{Spectral graph theory proper.} Cvetkovi\'c--Doob--Sachs
\cite{CDS95} is the classical monograph; its Chapter~1 contains the
coefficient theorem for digraphs with proof. Brouwer--Haemers
\cite{BH12} (freely available from the authors' web pages) is a modern,
readable account of the symmetric/undirected theory --- useful
background, with the caveat that most of it does not apply to the
signed, nonsymmetric matrices studied here. Saunders--Schneider
\cite{SS78} proves the equivalence between diagonal similarity and
equality of cycle products.

\section{The Python verification script}\label{sec:script}

The following script (i) builds $J_n(\overline{j})$ for a random
integer sequence with $j_1 = 1$; (ii) computes its characteristic
polynomial with \texttt{numpy}; (iii) recomputes the coefficients from
the digraph coefficient theorem, enumerating all sets of pairwise
disjoint intervals and using the cycle weights
\eqref{eq:cycleweight}; and (iv) compares both against the formula of
Theorem 2.3, with $h$ recovered from $j$ via equation (D1). All three
computations agree.

\begin{lstlisting}
import numpy as np
import random
from math import comb, factorial

random.seed(1)
n = 7
# random integer j-sequence with j_1 = 1 (as in Lemma 2.2)
j = [None, 1] + [random.randint(-4, 4) for _ in range(2, n + 1)]  # j[1..n]

# Build J_n(jbar): lower part Toeplitz j_{i-k+1}, superdiagonal -i
J = np.zeros((n, n))
for i in range(1, n + 1):
    for k in range(1, n + 1):
        if k <= i:
            J[i - 1, k - 1] = j[i - k + 1]
        elif k == i + 1:
            J[i - 1, k - 1] = -i

# Characteristic polynomial via numpy:
# coefficients of x^n + a_1 x^{n-1} + ...
cp_numeric = np.poly(J)

# --- digraph coefficient theorem ---
# Cycles of the digraph = intervals [k,m]:
#   weight = j_{m-k+1} * (-1)^(m-k) * (m-1)!/(k-1)!
def cycle_weight(k, m):
    return j[m - k + 1] * ((-1) ** (m - k)) * factorial(m - 1) // factorial(k - 1)

# Coefficient of x^{n-r}: sum over sets of disjoint intervals
# covering r vertices, with sign (-1)^{#cycles}  (det(xI - A) convention)
intervals = [(k, m) for k in range(1, n + 1) for m in range(k, n + 1)]

def disjoint_sets(ivs):
    """Enumerate all sets of pairwise disjoint intervals."""
    out = [[]]
    def rec(start, chosen):
        for idx in range(start, len(ivs)):
            iv = ivs[idx]
            if all(iv[0] > c[1] or iv[1] < c[0] for c in chosen):
                chosen.append(iv)
                out.append(list(chosen))
                rec(idx + 1, chosen)
                chosen.pop()
    rec(0, [])
    return out

coeffs = [0] * (n + 1)  # coeffs[r] multiplies x^{n-r}
coeffs[0] = 1
for S in disjoint_sets(intervals):
    if not S:
        continue
    r = sum(m - k + 1 for k, m in S)
    w = 1
    for k, m in S:
        w *= cycle_weight(k, m)
    coeffs[r] += ((-1) ** len(S)) * w

print("numpy char poly coeffs:    ", np.round(cp_numeric, 6))
print("digraph cycle-cover coeffs:", coeffs)
print("match:", np.allclose(cp_numeric, coeffs))

# --- confirm these equal Theorem 2.3's binom(n,r) r! h_r (-1)^r,
#     with h recovered from j via equation (D1) ---
h = [1]
for m in range(1, n + 1):
    s = sum(j[r] * h[m - r] for r in range(1, m + 1))
    h.append(s / m)
thm = [comb(n, r) * factorial(r) * h[r] * (-1) ** r for r in range(n + 1)]
print("Theorem 2.3 coeffs:        ", thm)
\end{lstlisting}

\noindent
Output for the seeded run ($n = 7$):
\begin{lstlisting}[language={}]
numpy char poly coeffs:     [ 1. -7. -21. 385. -805. -609. 6433. -8041.]
digraph cycle-cover coeffs: [1, -7, -21, 385, -805, -609, 6433, -8041]
match: True
Theorem 2.3 coeffs:         [1, -7.0, -21.0, 385.0, -805.0, -609.0,
                             6433.0, -8041.0]
\end{lstlisting}

\begin{thebibliography}{99}

\bibitem{BP94}
A.~Berman and R.~J. Plemmons,
\emph{Nonnegative Matrices in the Mathematical Sciences},
Classics in Applied Mathematics 9, SIAM, Philadelphia, 1994.
[Chapter~2 proves the Perron--Frobenius theorem and the equivalence of
irreducibility with strong connectivity of the digraph.]

\bibitem{BH12}
A.~E. Brouwer and W.~H. Haemers,
\emph{Spectra of Graphs},
Universitext, Springer, New York, 2012.
[A readable modern account of spectral graph theory for symmetric
(undirected) adjacency matrices, including Perron--Frobenius and
eigenvalue interlacing with proofs; freely available from the authors'
web pages.]

\bibitem{Brualdi82}
R.~A. Brualdi,
Matrices, eigenvalues, and directed graphs,
\emph{Linear and Multilinear Algebra} \textbf{11} (1982), 143--165.
[Original proof of the cycle-based refinement of the Ger\v{s}gorin disc
theorem used in these remarks; the digraph structure of the matrix
enters the proof directly.]

\bibitem{BC09}
R.~A. Brualdi and D.~Cvetkovi\'c,
\emph{A Combinatorial Approach to Matrix Theory and Its Applications},
Discrete Mathematics and Its Applications, CRC Press, Boca Raton, 2009.
[Develops determinants, characteristic polynomials, and the
Cayley--Hamilton theorem from scratch in the language of weighted
digraphs and cycle covers. The most accessible source for the
coefficient theorem; assumes only elementary linear algebra.]

\bibitem{CDS95}
D.~Cvetkovi\'c, M.~Doob, and H.~Sachs,
\emph{Spectra of Graphs: Theory and Applications},
3rd ed., Johann Ambrosius Barth, Heidelberg, 1995.
[The classical monograph on graph spectra. Chapter~1 states and proves
the coefficient theorem for the characteristic polynomial of a weighted
digraph, the form used in these remarks.]

\bibitem{Ger31}
S.~Ger\v{s}gorin,
\"Uber die Abgrenzung der Eigenwerte einer Matrix,
\emph{Izv. Akad. Nauk SSSR, Ser. Mat.} \textbf{7} (1931), 749--754.
[The original disc theorem; cited for history --- modern proofs are in
\cite{HJ13} and \cite{Varga04}.]

\bibitem{GVL13}
G.~H. Golub and C.~F. Van Loan,
\emph{Matrix Computations},
4th ed., Johns Hopkins University Press, Baltimore, 2013.
[Section 7.4 treats unreduced Hessenberg matrices, including the rank
argument showing every eigenvalue has geometric multiplicity one.]

\bibitem{HJ13}
R.~A. Horn and C.~R. Johnson,
\emph{Matrix Analysis},
2nd ed., Cambridge University Press, Cambridge, 2013.
[The standard reference. Chapter~6 proves the Ger\v{s}gorin disc
theorem, Taussky's boundary refinement, and the equivalence of
irreducibility with strong connectivity; \S 4.3 proves Cauchy
interlacing for Hermitian matrices; Chapter~8 proves Perron--Frobenius.
This is the same book cited as Horn and Johnson in the main paper.]

\bibitem{Meyer00}
C.~D. Meyer,
\emph{Matrix Analysis and Applied Linear Algebra},
SIAM, Philadelphia, 2000.
[Chapter~8 contains what is perhaps the most readable proof of the
Perron--Frobenius theorem, together with the digraph characterization
of irreducibility.]

\bibitem{SS78}
B.~D. Saunders and H.~Schneider,
Flows on graphs applied to diagonal similarity and diagonal equivalence
for matrices,
\emph{Discrete Math.} \textbf{24} (1978), 205--220.
[Proves that two matrices supported on the same digraph are diagonally
similar if and only if corresponding cycle products are equal --- the
precise sense in which the cycle weights \eqref{eq:cycleweight} are the
complete intrinsic spectral data.]

\bibitem{Stanley12}
R.~P. Stanley,
\emph{Enumerative Combinatorics, Volume 1},
2nd ed., Cambridge Studies in Advanced Mathematics 49,
Cambridge University Press, Cambridge, 2012.
[Section 4.7 (the transfer-matrix method) contains a short
self-contained proof of the cycle-cover expansion of $\det(I - xA)$,
equivalent to the coefficient theorem.]

\bibitem{Taussky49}
O.~Taussky,
A recurring theorem on determinants,
\emph{Amer. Math. Monthly} \textbf{56} (1949), 672--676.
[Source of the boundary refinement of Ger\v{s}gorin for irreducible
matrices.]

\bibitem{Varga04}
R.~S. Varga,
\emph{Ger\v{s}gorin and His Circles},
Springer Series in Computational Mathematics 36, Springer, Berlin, 2004.
[A short book devoted entirely to eigenvalue-inclusion theorems,
written to be self-contained. Chapter~1 proves the disc theorem and
Taussky's theorem; Chapter~2 proves Brualdi's cycle-based inclusion
sets. Recommended as the main reference for a reader new to this
material.]

\end{thebibliography}

\end{document}
